\section{An analytic period substitution and the row-sign correction}\label{app:period}

This appendix records a change of variables and a sign convention. Relations explicitly marked numerical refer to the historical candidate model and are not critical-value theorems.

\subsection{An exact change of variables}
Let
\[
 F(\lambda)={}_2F_1\!\left(\tfrac12,\tfrac12;1;\lambda\right),
 \qquad F_c(\lambda)=F(1-\lambda),\qquad 0<\lambda<1.
\]
For an integer $k\geq0$, an integer $0\leq j\leq k$, and a real number
$e>0$, define $y=F_c/F$. The inverse parameter is the modular lambda
function at $iy$. Put
\[
 G_e(iy)=\lambda(iy)^e\bigl(1-\lambda(iy)\bigr)
          \theta_3(0\mid iy)^{2k+4}.
\]
Then the following analytic identity holds:
\begin{equation}\label{eq:period-dictionary-static}
 \int_0^1\lambda^{e-1} F(\lambda)^j F_c(\lambda)^{k-j}\,d\lambda
 =\pi\int_0^\infty G_e(iy)y^{k-j}\,dy.
\end{equation}

\begin{proof}
Both $u=F$ and $v=F_c$ solve
\[
 \lambda(1-\lambda)u''+(1-2\lambda)u'-\tfrac14u=0.
\]
This is the specialization of the hypergeometric differential equation in
DLMF~15.10.1 \cite{DLMF}. Abel's identity therefore gives
\[
 uv'-u'v=\frac{C}{\lambda(1-\lambda)}.
\]
The hypergeometric--elliptic identity and the logarithmic endpoint expansion
(DLMF~19.5.1 and 19.12.1) yield
\[
 F(\lambda)=1+O(\lambda),\qquad
 F_c(\lambda)=\frac1\pi\log\frac{16}{\lambda}
                +O(\lambda|\log\lambda|)
 \quad(\lambda\longrightarrow0^+).
\]
The differentiated expansions imply $C=-1/\pi$. Hence
\begin{equation}\label{eq:period-jacobian-static}
 \frac{dy}{d\lambda}
 =-\frac{1}{\pi\lambda(1-\lambda)F(\lambda)^2}.
\end{equation}
The positive function $y$ decreases from $+\infty$ to $0$. DLMF~20.9.2,
together with DLMF~19.5.1, identifies
$F(\lambda(iy))=\theta_3(0\mid iy)^2$ and
$F_c(\lambda(iy))/F(\lambda(iy))=y$.
Using $F_c=yF$ and
$d\lambda=-\pi\lambda(1-\lambda)F^2\,dy$ proves
\eqref{eq:period-dictionary-static} after reversing the endpoints.
The original integrand is
$O(\lambda^{e-1}(1+|\log\lambda|)^{k-j})$ near $0$ and
$O((1+|\log(1-\lambda)|)^j)$ near $1$.
Both bounds are integrable under the stated hypotheses, so the improper
change of variables is justified.\end{proof}

For the present candidate periods, $k=4$ and
$e\in\{1/3,2/3,4/3,5/3\}$; writing $m=k+1-j$ gives the factor
$y^{m-1}$ used in the old note N12. The same argument also covers its
$k=2$ control integrals.
This establishes the analytic dictionary and its normalization. Identifying
its right-hand side with a specified arithmetic $L$-value requires further
modularity, rational-structure, and Mellin-normalization arguments. In
particular, the rational Scholl realization in Section~\ref{sec:scholl}
does not by itself identify these integrals with rational period bases or
prove the rational determinant constants below.

\subsection{The documented and implemented signs differ}
Write $H_j(\lambda)=F(\lambda)^jF_c(\lambda)^{4-j}$ and
\[
 C_c(j)=\int_0^1(\lambda^c+\lambda^{-c})H_j(\lambda)\,d\lambda.
\]
For two row families $A_{2/3},A_{1/3}$ use the fixed row order
\[
 \mu_A(j_1,j_2)=
 \det\begin{pmatrix}
 A_{2/3}(j_1)&A_{2/3}(j_2)\\
 A_{1/3}(j_1)&A_{1/3}(j_2)
 \end{pmatrix}.
\]
N12, line~45, describes the uniform convention
\[
 S_c^{\mathrm{unif}}(j)
 =\int_0^1(\lambda^c-\lambda^{-c})H_j(\lambda)\,d\lambda.
\]
The actual N11 implementation instead uses
\begin{align*}
 S_{2/3}^{\mathrm{code}}(j)
   &=\int_0^1(\lambda^{-2/3}-\lambda^{2/3})H_j(\lambda)\,d\lambda,\\
 S_{1/3}^{\mathrm{code}}(j)
   &=\int_0^1(\lambda^{1/3}-\lambda^{-1/3})H_j(\lambda)\,d\lambda.
\end{align*}
Thus, exactly and without a numerical computation,
\[
 \mu_{S^{\mathrm{code}}}(j_1,j_2)
 =-\mu_{S^{\mathrm{unif}}}(j_1,j_2).
\]
The old negative ratios use the implemented convention. To retain the
uniform convention in prose, the corresponding \emph{reported numerical}
relations must instead read
\begin{align*}
 \frac{\mu_{S^{\mathrm{unif}}}(1,3)}{\mu_C(2,4)}
     &\mathrel{\stackrel{\mathrm{num}}{=}}\frac12, &
 \frac{\mu_{S^{\mathrm{unif}}}(0,2)}{\mu_C(1,3)}
     &\mathrel{\stackrel{\mathrm{num}}{=}}\frac43,\\
 \mathcal L_{\rm candidate}(4)&\mathrel{\stackrel{\mathrm{num}}{=}}
       \frac{\pi^6\mu_{S^{\mathrm{unif}}}(0,2)}{209952}.
\end{align*}
Alternatively, explicitly define the two implemented row orientations and
retain the negative signs in the historical outputs. No numerical rerun is
needed to diagnose this discrepancy. The five principal critical-value
relations expressed only through $\mu_C$ are unaffected.

